%% kapa: exam=1 round=25 year=2014 period=1 form=A subject=civil_law subject_ko=민법 seq=3 number=3
%% source: Q-Net 감정평가사 기출문제 (cst00309) artlSeq=5101207 file=extracted/1008/5101207/2014년 제25회 감정평가사 1차시험 기출문제/1교시(A형).hwp
%% conversion: kordoc 4.24.51 -> tools/gen1.py (auto); review=proofread-sonnet; flags=none
\begin{question}{3}
제한능력제도에 관한 설명으로 옳지 않은 것은?
\begin{choices}
  \item 피성년후견인이 행한 모든 재산법상의 법률행위는 언제든지 취소할 수 있다.
  \item 가정법원은 한정후견개시의 심판을 할 때 본인의 의사를 고려하여야 한다.
  \item 가정법원은 특정후견의 심판을 할 때 본인의 의사에 반하여 할 수 없다.
  \item 가정법원이 피한정후견인에 대해 성년후견 개시의 심판을 하는 경우, 종전의 한정후견의 종료심판을 한다.
  \item 가정법원은 청구권자의 청구가 없는 한, 성년후견개시의 심판을 직권으로 할 수 없다.
\end{choices}
\end{question}
\answer{1}
